\chapter{2011年上海卷（秋理）}

\section{填空题}

\begin{problem}
函数 \( f(x) = \frac{1}{x - 2} \) 的反函数为 \( f^{-1}(x) = \)\fillinblank{}.
\end{problem}

\begin{problem}
若全集 \(U = \R\)，集合 \(A = \{x \mid x \geqslant 1\} \cup \{x \mid x \leqslant 0\}\)，则 \(\complement_U A =\fillinblank{}\).
\end{problem}

\begin{problem}
设 \(m\) 为常数，若点 \(F(0,5)\) 是双曲线 \(\frac{y^2}{m} - \frac{x^2}{9} = 1\) 的一个焦点，则 \(m =\fillinblank{}\).
\end{problem}

\begin{problem}
不等式 \(\frac{x + 1}{x} \leqslant 3\) 的解集为\fillinblank{}.
\end{problem}

\begin{problem}
在极坐标系中，直线 \(\rho(2\cos \theta + \sin \theta) = 2\) 与直线 \(\rho\cos \theta = 1\) 的夹角大小为\fillinblank{}. （结果用反三角函数值表示）
\end{problem}

\begin{problem}
在相距 2 千米的 A, B 两点处测量目标 C，若 \(\angle CAB = 75^{\circ}\)，\(\angle CBA = 60^{\circ}\)，则 A, C 两点之间的距离为\fillinblank{}\text{千米}.
\end{problem}

\begin{problem}
若圆锥的侧面积为 \(2\pi\)，底面积为 \(\pi\)，则该圆锥的体积为\fillinblank{}.
\end{problem}

\begin{problem}
函数 \( y = \sin \left(\frac{\pi}{2} + x\right) \cos \left(\frac{\pi}{6} - x\right) \) 的最大值为\fillinblank{}.
\end{problem}

\begin{problem}
马老师从课本上抄录一个随机变量 \(\xi\) 的概率分布列如下表：

\begin{center}
\begin{tabular}{|c|c|c|c|}
\hline
\(X\) & 1 & 2 & 3 \\
\hline
\(P(\xi=x)\) & ? & ! & ? \\
\hline
\end{tabular}
\end{center}

请小牛同学计算 \( \xi \) 的数学期望，尽管“!”处无法完全看清，且两个“?”处字迹模糊，但能肯定这两个“?”处的数值相同. 据此，小牛给出了正确答案 \(E(\xi)=\)\fillinblank{}.
\end{problem}

\begin{problem}
行列式 \(\left| \begin{array}{cc}a & b\\ c & d \end{array} \right|\) \((a,b,c,d\in \{-1,1,2\})\) 所有可能的值中，最大的是\fillinblank{}.
\end{problem}

\begin{problem}
在正三角形 \(ABC\) 中，\(D\) 是 \(BC\) 上的点. 若 \(AB = 3\)，\(BD = 1\)，则 \(\overrightarrow{AB} \cdot \overrightarrow{AD} =\fillinblank{}.\)
\end{problem}

\begin{problem}
随机抽取的 9 位同学中，至少有 2 位同学在同一月份出生的概率为\fillinblank{}. （默认每个月的天数相同，结果精确到 0.001）
\end{problem}

\begin{problem}
设 \(g(x)\) 是定义在 \(\R\) 上，以 1 为周期的函数，若函数 \(f(x) = x + g(x)\) 在区间 \([3,4]\) 上的值域为 \([-2,5]\)，则 \(f(x)\) 在区间 \([-10,10]\) 上的值域为\fillinblank{}.
\end{problem}

\begin{problem}
已知点 \(O(0,0)\)，\(Q_0(0,1)\) 和点 \(R_0(3,1)\)，记 \(Q_0R_0\) 的中点为 \(P_1\)，取 \(Q_0P_1\) 和 \(P_1R_0\) 中的一条，记其端点为 \(Q_1\)，\(R_1\)，使之满足 \((|OQ_1| - 2)(|OR_1| - 2) < 0\)，记 \(Q_1R_1\) 的中点为 \(P_2\)，取 \(Q_1P_2\) 和 \(P_2R_1\) 中的一条，记其端点为 \(Q_2\)，\(R_2\)，使之满足 \((|OQ_2| - 2)(|OR_2| - 2) < 0\). 依次下去，得到 \(P_1\)，\(P_2\)，\(\cdots\)，\(P_n\)，\(\cdots\)，则 \(\lim\limits_{n \to \infty} \left|Q_0P_n\right| =\fillinblank{} \).
\end{problem}

\section{单选题}

\begin{problem}
若 \(a\)，\(b \in \R\)，且 \(ab > 0\)，则下列不等式中，恒成立的是（\quad）

\choices
    {\(a^{2} + b^{2} > 2ab\)}
    {\(a + b \geqslant 2\sqrt{ab}\)}
    {\(\frac{1}{a} + \frac{1}{b} > \frac{2}{\sqrt{ab}}\)}
    {\(\frac{b}{a} + \frac{a}{b} \geqslant 2\)}
\end{problem}

\begin{problem}
下列函数中，既是偶函数，又是在区间 \((0,+\infty)\) 上单调递减的函数是（\quad）

\choices
    {\(y = \ln \frac{1}{|x|}\)}
    {\(y = x^{3}\)}
    {\(y = 2^{|x|}\)}
    {\(y = \cos x\)}
\end{problem}

\begin{problem}
设 \(A_1\)，\(A_2\)，\(A_3\)，\(A_4\)，\(A_5\) 是平面上给定的 5 个不同点，则使 \(\overrightarrow{MA_1} + \overrightarrow{MA_2} + \overrightarrow{MA_3} + \overrightarrow{MA_4} + \overrightarrow{MA_5} = \vec{0}\) 成立的点 \(M\) 的个数为（\quad）

\choices
    {\(0\)}
    {\(1\)}
    {\(5\)}
    {\(10\)}
\end{problem}

\begin{problem}
设 \(\{a_n\}\) 是各项为正数的无穷数列，\(A_i\) 是边长为 \(a_i\)，\(a_{i+1}\) 的矩形面积 \((i = 1, 2, \cdots)\)，则 \(\{A_n\}\) 为等比数列的充要条件为（\quad）

\choices
    {\(\{a_n\}\) 是等比数列}
    {\(\{a_1, a_3, \cdots, a_{2n-1}, \cdots\}\) 或 \(\{a_2, a_4, \cdots, a_{2n}, \cdots\}\) 是等比数列}
    {\(\{a_1, a_3, \cdots, a_{2n-1}, \cdots\}\) 和 \(\{a_2, a_4, \cdots, a_{2n}, \cdots\}\) 均是等比数列}
    {\(\{a_1, a_3, \cdots, a_{2n-1}, \cdots\}\) 和 \(\{a_2, a_4, \cdots, a_{2n}, \cdots\}\) 均是等比数列，且公比相同}
\end{problem}

\section{解答题}

\begin{problem}
已知复数 \(z_1\) 满足 \((z_1 - 2)(1 + i) = 1 - i\)（\(i\) 为虚数单位），复数 \(z_2\) 的虚部为 2，且 \(z_1 \cdot z_2\) 是实数，求 \(z_2\).
\end{problem}

\begin{problem}
已知函数 \(f(x) = a \cdot 2^{x} + b \cdot 3^{x}\)，其中常数 \(a\)，\(b\) 满足 \(a \cdot b \neq 0\).

\begin{enumerate}
\item 若 \(a \cdot b > 0\)，判断函数 \(f(x)\) 的单调性；
\item 若 \(a \cdot b < 0\)，求 \(f(x+1) > f(x)\) 时的 \(x\) 的取值范围.
\end{enumerate}
\end{problem}

\begin{problem}
已知 \(ABCD - A_{1}B_{1}C_{1}D_{1}\) 是底面边长为 1 的正四棱柱，\(O_1\) 为 \(A_1C_1\) 与 \(B_1D_1\) 的交点.

\begin{enumerate}
\item 设 \(AB_1\) 与底面 \(A_1B_1C_1D_1\) 所成角的大小为 \(\alpha\)，二面角 \(A - B_{1}D_{1} - A_1\) 的大小为 \(\beta\). 求证：\(\tan \beta = \sqrt{2}\tan \alpha\)；
\item 若点 \(C\) 到平面 \(AB_1D_1\) 的距离为 \(\frac{4}{3}\)，求正四棱柱 \(ABCD - A_{1}B_{1}C_{1}D_{1}\) 的高.
\end{enumerate}

\begin{center}
\bitmapinclude[max width=0.34\linewidth,max totalheight=4.8cm]{img/2011/shanghai_science/q21_fig1.png}
\end{center}
\end{problem}

\begin{problem}
已知数列 \(\{a_n\}\) 和 \(\{b_n\}\) 的通项公式分别为 \(a_n = 3n + 6\)，\(b_n = 2n + 7\) \((n \in \mathbf{N}^{*})\)，将集合 \(\{x \mid x = a_n, n \in \mathbf{N}^{*}\} \cup \{x \mid x = b_n, n \in \mathbf{N}^{*}\}\) 中的元素从小到大依次排列，构成数列 \(c_1\)，\(c_2\)，\(c_3\)，\(\cdots\)，\(c_n\)，\(\cdots\).

\begin{enumerate}
\item 求 \(c_1\)，\(c_2\)，\(c_3\)，\(c_4\)；
\item 求证：在数列 \(\{a_n\}\) 中，但不在数列 \(\{b_n\}\) 中的项恰为 \(a_2\)，\(a_4\)，\(\cdots\)，\(a_{2n}\)，\(\cdots\)；
\item 求数列 \(\{c_n\}\) 的通项公式.
\end{enumerate}
\end{problem}

\begin{problem}
已知平面上的线段 \(l\) 及点 \(P\)，在 \(l\) 上任取一点 \(Q\)，线段 \(PQ\) 长度的最小值称为点 \(P\) 到线段 \(l\) 的距离，记作 \(d(P, l)\).

\begin{enumerate}
\item 求点 \(P(1,1)\) 到线段 \(l: x - y - 3 = 0 (3 \leqslant x \leqslant 5)\) 的距离 \(d(P, l)\)；
\item 设 \(l\) 是长为 2 的线段，求点的集合 \(D = \{P \mid d(P, l) \leqslant 1\}\) 所表示的图形面积；
\item 写出到两条线段 \(l_1\)，\(l_2\) 距离相等的点的集合 \(\Omega = \{P \mid d(P, l_1) = d(P, l_2)\}\)，其中 \(l_1 = AB\)，\(l_2 = CD\)，\(A\)，\(B\)，\(C\)，\(D\) 是下列三组点中的一组.

以下三组点任选一组：

\begin{circlelist}
\item \(A(1,3)\)，\(B(1,0)\)，\(C(-1,3)\)，\(D(-1,0)\).
\item \(A(1,3)\)，\(B(1,0)\)，\(C(-1,3)\)，\(D(-1,-2)\).
\item \(A(0,1)\)，\(B(0,0)\)，\(C(0,0)\)，\(D(2,0)\).
\end{circlelist}
\end{enumerate}
\end{problem}
